\documentclass[11pt,a4paper]{article}
\usepackage[margin=1in]{geometry}
\usepackage{amsmath}
\usepackage{amssymb}
\usepackage{physics}
\usepackage{xcolor}
\usepackage{listings}
\usepackage{hyperref}
\usepackage{graphicx}
\usepackage{booktabs}
\usepackage{array}
\usepackage{fancybox}

% Configure listings for Python code
\lstset{
    language=Python,
    basicstyle=\ttfamily\small,
    breaklines=true,
    columns=fullflexible,
    keywordstyle=\color{blue},
    commentstyle=\color{gray},
    stringstyle=\color{orange},
    showstringspaces=false,
    frame=tb
}

\title{Signed SMEFT Weight Decomposition into Pos/Neg Flows}
\author{MLHEPsim Flow Training Pipeline}
\date{\today}

\begin{document}

\maketitle

\begin{abstract}
This document describes the current implementation of normalising flow training for SMEFT event reweighting.
The core innovation is the decomposition of signed weight components (linear and quadratic SMEFT terms) into 
non-negative positive and negative parts, enabling stable training of five independent normalising flows.
This replaces previous attempts at clamping, absolute value corrections, and per-batch rescaling with an 
exact and principled solution to the fundamental incompatibility between normalising flows (which require 
non-negative weights) and signed SMEFT measures.
\end{abstract}

\section{Introduction}

\subsection{The Problem: Signed Weights in Normalising Flows}

Standard normalising flow training assumes non-negative importance weights. However, the SMEFT decomposition 
of the top-pair production cross section produces intrinsically signed weight components:

\begin{equation}
w_{\text{full}}(x; c_{\text{Ht}}) = w_{\text{sm}}(x) + c_{\text{Ht}} \, w_{\text{lin}}(x) + c_{\text{Ht}}^2 \, w_{\text{quad}}(x)
\end{equation}

where:
\begin{itemize}
  \item $w_{\text{sm}}(x)$: SM weight (always positive)
  \item $w_{\text{lin}}(x)$: Linear interference term (\textbf{$\approx 16.8\%$ negative events})
  \item $w_{\text{quad}}(x)$: Quadratic term (\textbf{$\approx 0.04\%$ negative events — essentially negligible})
\end{itemize}

\noindent Statistics from 3.7M events (from \texttt{ttZ\_for\_Melle\_tree.root}):
\begin{table}[h!]
\centering
\begin{tabular}{lccccc}
\toprule
Component & Events & \% Total & Z-Factor & \% Negative & Role \\
\midrule
Linear (+) & 3M & 83.2\% & 54.568 & — & Primary \\[1pt]
Linear (−) & 624k & 16.8\% & 4.542 & ✓ & Sign-split \\[1pt]
Quadratic (+) & 3.7M & 99.96\% & 3.014 & — & Primary \\[1pt]
Quadratic (−) & 1.5k & 0.04\% & 0.00465 & ✓ & Negligible \\[1pt]
\bottomrule
\end{tabular}
\end{table}

When training a normalising flow $q(x)$ on such data, the objective becomes sign-indefinite and unbounded:
\begin{equation}
\mathcal{L} = \mathbb{E}_{x \sim \text{data}}[w(x) \cdot (-\log q(x))]
\end{equation}

For a negative-weight event with $w_i < 0$, once the flow learns and $-\log q(x_i)$ becomes large, 
the product $w_i \cdot (-\log q(x_i))$ becomes a large \emph{positive} loss, destabilising gradient descent.

\subsection{Why Workarounds Don't Work}

Previous attempts included clamping, absolute values, and per-batch SNIS normalisation. All modify what 
the flows learn, distorting the physics:
\begin{itemize}
  \item \textbf{Clamping to zero}: Discards physically meaningful negative amplitudes
  \item \textbf{Taking absolute values}: Breaks the signed interference pattern
  \item \textbf{SNIS rescaling}: Prevents overflow but doesn't address that the objective is sign-indefinite
\end{itemize}

\subsection{The Solution: Pos/Neg Decomposition}

Train \textbf{five independent flows}, one per non-negative weight component:

\begin{equation}
\begin{aligned}
w_{\text{lin}}(x) &= w_{\text{lin}}^+(x) - w_{\text{lin}}^-(x) \\
w_{\text{quad}}(x) &= w_{\text{quad}}^+(x) - w_{\text{quad}}^-(x)
\end{aligned}
\end{equation}

where:
\begin{equation}
\begin{aligned}
w_k^+(x) &:= \max(w_k(x), 0) \\
w_k^-(x) &:= \max(-w_k(x), 0)
\end{aligned}
\end{equation}

Each positive/negative part is a valid non-negative importance weight, so the objective is well-posed. 
The combination formula at inference reconstructs the exact signed SMEFT measure.

---

\section{The Five Flows}

\subsection{Flow Training}

\begin{table}[h!]
\centering
\begin{tabular}{llcc}
\toprule
Flow Name & Trained On & Training Weight & Events Used \\
\midrule
\texttt{sm} & All events & $w_{\text{sm}}(x)$ & $N$ \\
\texttt{linear\_pos} & $w_{\text{lin}}(x) \geq 0$ & $w_{\text{lin}}(x)$ & $\sim 83.2\% \, N$ \\
\texttt{linear\_neg} & $w_{\text{lin}}(x) < 0$ & $|w_{\text{lin}}(x)|$ & $\sim 16.8\% \, N$ \\
\texttt{quadratic\_pos} & $w_{\text{quad}}(x) \geq 0$ & $w_{\text{quad}}(x)$ & $\sim 99.96\% \, N$ \\
\texttt{quadratic\_neg} & $w_{\text{quad}}(x) < 0$ & $|w_{\text{quad}}(x)|$ & $\sim 0.04\% \, N$ (sparse, tiny weights) \\
\bottomrule
\end{tabular}
\end{table}

\noindent Each flow learns:
\begin{equation}
q_k(x) \approx \frac{|w_k(x)| \cdot p_0(x)}{\int |w_k(x')| \, p_0(x') \, dx'} = \frac{|w_k(x)| \, p_0(x)}{Z_k}
\end{equation}

where $p_0(x)$ is the true SM distribution and $Z_k := \sum_i |w_{k,i}|$ is the normalisation constant.

\subsection{Key Properties}

\begin{itemize}
  \item \textbf{All training weights are strictly non-negative}—objective is well-posed
  \item \textbf{No information is discarded}—all events contribute to one of the five flows
  \item \textbf{Exact reconstruction}—the combination formula reproduces the original signed measure
  \item \textbf{Signed weights at inference}—negative weights can still appear (destructive interference)
\end{itemize}

---

\section{Implementation: Data Loading}

\subsection{Dataset Setup in \texttt{ttz\_dataset.py}}

\subsubsection{Weight Column Mapping}

\begin{lstlisting}
WEIGHT_COLS = {
    "sm": 15, "linear": 16, "quadratic": 17, "full": 18,
    "linear_pos": 16, "linear_neg": 16,
    "quadratic_pos": 17, "quadratic_neg": 17,
}
\end{lstlisting}

The \texttt{\_pos} and \texttt{\_neg} variants point to the same raw column but apply sign-based filtering in \texttt{setup()}.

\subsubsection{Sign-Based Event Filtering}

In the \texttt{ttzDataModule.setup()} method:

\begin{lstlisting}
# Example: linear_pos type
if self.weight_type.endswith('_pos'):
    valid_mask = self.weights >= 0
    data = data[valid_mask]
    self.weights = self.weights[valid_mask]
    # Mean normalisation computed only on kept events
    self.weights = self.weights / np.mean(self.weights)

# Example: linear_neg type
elif self.weight_type.endswith('_neg'):
    valid_mask = self.weights < 0
    data = data[valid_mask]
    self.weights = -self.weights[valid_mask]  # Flip sign -> positive
    self.weights = self.weights / np.mean(self.weights)
\end{lstlisting}

\noindent Key points:
\begin{enumerate}
  \item Filter by sign \textbf{before} mean normalisation
  \item For \texttt{\_neg}, flip sign: $|w| > 0$
  \item Compute $\text{mean}(w)$ only on kept subset
  \item All training weights $\to$ positive
\end{enumerate}

\subsection{Stratified Sampling}

Both train and validation dataloaders use \texttt{WeightStratifiedSampler}, which divides events into 
4 quantile bins by $|w|$ and constructs each batch by drawing equally from all bins. This prevents 
batches dominated by a handful of extreme-weight events.

---

\section{Normalisation Constants}

\subsection{Computation from ROOT Data}

In \texttt{smeft\_reweighting.py}:

\begin{lstlisting}
def compute_z_normalisations(root_file_path: str) -> dict:
    """Compute Z_k = sum_i w_k(i) for each weight component."""
    
    with uproot.open(root_file_path) as f:
        tree = f["Events"]
        # Load using awkward to handle ragged arrays correctly
        data = tree.arrays(['eventWeight', 'smeft_weights'], 
                          library='ak')
    
    w_sm = ak.to_numpy(data['eventWeight'])
    w_plus  = ak.to_numpy(data['smeft_weights'][:, IDX_CHT_PLUS5])
    w_minus = ak.to_numpy(data['smeft_weights'][:, IDX_CHT_MINUS5])
    
    # Decompose (c = 5.0, c^2 = 25.0)
    w_lin  = (w_plus - w_minus) / 10.0
    w_quad = (w_plus + w_minus - 2.0 * w_sm) / 50.0
    
    # Compute Z values
    Z = {
        "sm":            float(w_sm.sum()),
        "linear":        float(w_lin.sum()),
        "linear_pos":    float(w_lin[w_lin > 0].sum()),
        "linear_neg":    float((-w_lin[w_lin < 0]).sum()),
        "quadratic":     float(w_quad.sum()),
        "quadratic_pos": float(w_quad[w_quad > 0].sum()),
        "quadratic_neg": float((-w_quad[w_quad < 0]).sum()),
    }
    
    # Sanity checks
    assert abs(Z["linear"] - (Z["linear_pos"] - Z["linear_neg"])) < 1e-3
    assert abs(Z["quadratic"] - (Z["quadratic_pos"] - Z["quadratic_neg"])) < 1e-3
    return Z
\end{lstlisting}

\subsection{Decomposition Verification}

The assertions verify that the decomposition is mathematically exact:

\begin{equation}
\begin{aligned}
Z_{\text{lin}} &= Z_{\text{lin}}^+ - Z_{\text{lin}}^- \\
Z_{\text{quad}} &= Z_{\text{quad}}^+ - Z_{\text{quad}}^-
\end{aligned}
\end{equation}

This ensures the five flows and combination formula can perfectly reconstruct the original signed distribution 
(up to numerical precision).

\subsection{Data Source Consistency}

All components of the pipeline use the same ROOT file source:
\begin{itemize}
  \item \textbf{Training data}: Generated by \texttt{process\_ttz\_dataset.py} from \\
        \texttt{/project/atlas/users/kdevries/EventLoop/ttZ\_for\_Melle\_tree.root} \\
        $\to$ produces \texttt{ttz\_weights.npy} (3.7M events, 19 columns)
  \item \textbf{Z-factor computation}: \texttt{smeft\_reweighting.py} reads the same file \\
        $\to$ ensures consistency with what flows were trained on
  \item \textbf{Weight contributions}: \texttt{plot\_weight\_contributions.py} also reads the same file \\
        $\to$ validates decomposition structure
\end{itemize}

\textbf{Important}: A different ROOT file (\texttt{full\_SR\_ttZ\_tree.root}) exists but is \textbf{not used} 
for training/reweighting. Mixing files causes Z-factor mismatches (e.g., $Z_{\text{quad}}^- = 0$ if the 
file has no negative quadratic events, when in fact $Z_{\text{quad}}^- = 0.00465$ in the training data).

---

\section{Density Ratios}

\subsection{Motivation: Why Density Ratios?}

The goal at inference time is to produce importance weights that reweight SM-generated events to look
like events from a SMEFT distribution at some Wilson coefficient $c_{\text{Ht}}$.  To see why density
ratios are the right tool, recall that each flow $k$ has learned the normalised distribution:

\begin{equation}
  q_k(x) \approx \frac{|w_k(x)|\,p_0(x)}{Z_k}
\end{equation}

where $p_0(x)$ is the true SM distribution and $Z_k = \int |w_k(x')|\,p_0(x')\,dx'$ is the
normalisation constant computed in Section~3.  Rearranging:

\begin{equation}
  |w_k(x)| = Z_k \cdot \frac{q_k(x)}{p_0(x)}
\end{equation}

The ratio $q_k(x)/p_0(x)$ is precisely the \textbf{density ratio} between flow $k$ and the SM flow.
We never need $q_k$ or $p_0$ individually in absolute terms---only their ratio matters.

On a set of events $\{x_i\}$ sampled from the SM flow, the density ratio becomes:

\begin{equation}
  r_k(x_i) = \frac{q_k(x_i)}{q_{\text{sm}}(x_i)}
\end{equation}

and the original weight is recovered as $|w_k(x_i)| \approx Z_k \cdot r_k(x_i)$.

\paragraph{Limitations of the approximation $q_{\text{sm}} \approx p_0$.}

It is worth being honest about what this approximation does and does not guarantee.  The density 
ratio $r_k(x) = q_k(x)/q_{\text{sm}}(x)$ recovers $|w_k(x)|/Z_k$ \emph{exactly} if and only if 
both $q_k$ and $q_{\text{sm}}$ are perfect.  In practice, neither is---they are neural networks 
trained on finite data with stochastic optimisation, and their errors are in general unpredictable 
and independent of each other.

One might hope that because all five flows are trained on the same underlying data, their 
approximation errors are correlated and therefore cancel in the ratio.  This is not guaranteed.  
Two networks with different weight columns, different sign-filtered training sets, different random 
initialisations, and different effective sample sizes will find different local minima and generalise 
differently.  There is no mechanism that reliably aligns their errors.

What actually justifies the approach is more modest:

\begin{enumerate}
  \item \textbf{Well-sampled regions are accurate individually.}  Both the SM flow and the auxiliary 
        flows receive many training examples in the bulk of phase space (e.g., moderate $p_T^Z$).  
        With enough data, a well-regularised normalising flow converges close to the true density.  
        In these regions, both $q_k$ and $q_{\text{sm}}$ should be accurate, so their ratio is 
        accurate---not because errors cancel, but because both errors are small.

  \item \textbf{Sampling from $q_{\text{sm}}$ concentrates evaluation in regions the SM flow knows 
        well.}  At inference time we draw $\{x_i\}$ from the SM flow itself.  The SM flow assigns 
        high probability to events it learned well, so the sampled events tend to live in regions 
        where $q_{\text{sm}}$ is reliable.  We are therefore unlikely to evaluate the ratio in 
        regions where $q_{\text{sm}}$ is wildly wrong---it would simply not sample there.

  \item \textbf{Sparse components carry a real bias risk.}  For components trained on far fewer events 
        (e.g., \texttt{linear\_neg} on 16.8\% of data, \texttt{quadratic\_neg} on only 0.04\%), the 
        auxiliary flow $q_k$ may be inaccurate in regions that are well-sampled by the SM flow.  
        The recovered weight is then biased---not by a small, controlled amount, but potentially 
        significantly.  This is a genuine limitation of the method.

  \item \textbf{The only real guarantee is empirical.}  If the density ratios are accurate, the 
        reweighted distributions produced by the combination formula should match the ground-truth 
        ROOT file distributions at each tested $c_{\text{Ht}}$ value (Section~\ref{sec:validation}).  
        Agreement in those comparison plots is the actual evidence that the approximation is 
        acceptable---not any theoretical argument about error structure.
\end{enumerate}

In summary: the ratio approach is a practical, well-motivated strategy, not a mathematically 
guaranteed one.  Its validity must be confirmed empirically for each flow component.

\subsection{Sampling Events from the SM Flow}

Before any density ratios can be computed, $N$ events must be drawn from the SM flow.  These events
$\{x_i\}_{i=1}^N$ serve as Monte Carlo integration points throughout the rest of the pipeline.

The SM flow $q_{\text{sm}}$ was trained on SM-weighted events from the ROOT file.  The ROOT file 
data inherently includes event selection cuts and detector effects from the simulation pipeline, 
so $q_{\text{sm}}$ learns the distribution of the \emph{selected, reconstructed} data, not the true 
underlying SM distribution.  Sampling from the trained flow therefore produces events with the same 
selection and reconstruction structure as the training data---the same biases are present.

This is actually desirable.  Because all five flows (SM, linear\_pos, linear\_neg, quadratic\_pos, 
quadratic\_neg) are trained on the \emph{same} ROOT file with the \emph{same} selection and detector 
effects, they all operate in the same biased space.  When the combination formula computes density 
ratios to recover physical weights, those ratios are formed between flows that share a common 
bias landscape.  This mutual coherence is more important than achieving an ideally unbiased sample.

In the code, \texttt{sample\_from\_sm\_flow()} draws samples in the preprocessed (scaled) space of
the flow.  After later applying the inverse scaler (\texttt{inverse\_transform}), the events are
returned to physical units (GeV, radians, etc.) suitable for plotting.  The resulting samples carry 
forward the phase-space structure learned from the ROOT file, including any selection artifacts.

\subsection{Log Probability Computation}

\label{sec:logp}

\subsubsection{Latent Space and Latent Codes}

Every normalising flow $q_k$ is defined by a bijective transformation $f_k: \mathcal{Z} \to \mathcal{X}$ that 
maps from a simple base distribution (a Mixture of Gaussians, MoG) \emph{in latent space} $\mathcal{Z}$ to the 
physical \emph{data space} $\mathcal{X}$ (our kinematic features).  

The \textbf{latent code} $z$ is the internal, learned compression of an event $x$---it is a 
hidden-layer representation obtained by inverting the forward flow: $z = f_k^{-1}(x)$.  Think of 
$z$ as the neural network's compact ``summary'' of the event $x$, living in an abstract space 
of dimension $D$ (typically 10--30 for this pipeline).  

\textbf{Note on notation:} The latent code $z$ (lowercase) and latent space $\mathcal{Z}$ (calligraphic) 
are \emph{not related} to the normalisation constants $Z_k$ (uppercase Z with subscript $k$) from 
Section~3.  The $Z_k$ are scalar sums ($Z_k = \sum_i |w_{k,i}|$, units: \emph{number of events}), 
whereas $z$ is a vector in abstract learned space (no physical units).  The similarity in notation 
can be confusing; keep these distinct concepts separate.

\subsubsection{Change-of-Variables Formula}

By the change-of-variables formula, the log-probability of a data point $x$ under flow $k$ is:

\begin{equation}
  \log q_k(x) = \log p_{\mathcal{Z}}\!\bigl(f_k^{-1}(x)\bigr)
                + \log\left|\det\frac{\partial f_k^{-1}}{\partial x}\right|
\end{equation}

The two terms are:
\begin{itemize}
  \item $\log p_{\mathcal{Z}}(z)$: log-probability of the latent code $z = f_k^{-1}(x)$ under the
        MoG base distribution in latent space.  This term reflects whether the latent representation 
        is typical for the learned base distribution---this is what \texttt{mog\_nll} computes (negated, 
        since NLL is returned).
  \item $\log|\det \partial f_k^{-1}/\partial x|$: log absolute Jacobian determinant of the inverse
        transformation.  This term accounts for the local volume dilation/contraction when mapping from 
        physical data space to latent space.  For MADE/MAF-type flows this is inexpensive to compute 
        since the Jacobian is triangular by construction.
\end{itemize}

In practice, \texttt{MADEMOG.estimate\_density(exp=False, mean=False)} returns the \emph{negative}
log-likelihood (NLL) per event, so we negate to obtain $\log q_k(x_i)$:

\begin{lstlisting}
nll = inner.estimate_density(x_scaled, exp=False, mean=False)  # (N,)
log_p = -nll   # NLL -> log-probability
\end{lstlisting}

This is evaluated for all five flows on the \emph{same} set of $N$ sampled events, giving five
arrays $\log q_{\text{sm}}, \log q_{\text{lin}^+}, \log q_{\text{lin}^-},
\log q_{\text{quad}^+}, \log q_{\text{quad}^-}$, each of length $N$.

\subsection{Log Density Ratios}

\label{sec:logr}

The log density ratio for each auxiliary flow relative to the SM flow is computed by straightforward
subtraction in log-space:

\begin{equation}
\begin{aligned}
\log r_{\text{lin},i}^{+}   &= \log q_{\text{lin}^+}(x_i) - \log q_{\text{sm}}(x_i) \\
\log r_{\text{lin},i}^{-}   &= \log q_{\text{lin}^-}(x_i) - \log q_{\text{sm}}(x_i) \\
\log r_{\text{quad},i}^{+}  &= \log q_{\text{quad}^+}(x_i) - \log q_{\text{sm}}(x_i) \\
\log r_{\text{quad},i}^{-}  &= \log q_{\text{quad}^-}(x_i) - \log q_{\text{sm}}(x_i)
\end{aligned}
\end{equation}

\paragraph{Physical interpretation.}
Each log ratio encodes \emph{how much more (or less) likely} a given event $x_i$ is under the
auxiliary weight distribution compared to the SM.

\begin{itemize}
  \item $\log r_k(x_i) > 0$ ($r_k > 1$): event $x_i$ is more probable under flow $k$ than under
        the SM---it lives in a region of phase space that is enhanced by the corresponding weight
        component.
  \item $\log r_k(x_i) < 0$ ($r_k < 1$): event $x_i$ is less probable under flow $k$---it is
        suppressed relative to SM in that region of phase space.
  \item $\log r_k(x_i) = 0$ ($r_k = 1$): flow $k$ is indistinguishable from the SM at this point;
        the weight component adds no information there.
\end{itemize}

For example, if $\log r_{\text{lin},i}^{+}$ is large and positive for events with high $p_T^Z$,
this tells us that the linear positive weight component is concentrated at high $p_T^Z$---consistent
with the EFT enhancing the high-energy tail of the distribution.

\paragraph{Why log-subtraction is numerically stable.}
Evaluating $q_k(x)/q_{\text{sm}}(x)$ directly by dividing probabilities would catastrophically
amplify rounding errors when probabilities are near zero.  Subtracting log-probabilities is exact up
to floating-point precision and avoids both underflow (near-zero probabilities) and overflow
(very high-probability events).

\subsection{Caching and Reuse}

Computing the log-probabilities for $N = 10^7$ events under five flows involves five separate
forward passes through the network---by far the most expensive step of the pipeline.  Consequently,
the four log-ratio arrays are saved immediately after computation:

\begin{lstlisting}
np.savez('density_ratios.npz',
         log_r_lin_pos  = log_r_lin_pos,
         log_r_lin_neg  = log_r_lin_neg,
         log_r_quad_pos = log_r_quad_pos,
         log_r_quad_neg = log_r_quad_neg)
\end{lstlisting}

On subsequent runs, \texttt{--load-ratios} skips the forward passes entirely and reads the saved
arrays directly.  Since evaluating the combination formula (Section~5) for a new $c_{\text{Ht}}$
value takes only microseconds (a handful of vector operations on $N$ numbers), the cached ratios
allow rapid exploration of many $c$ values at negligible cost.

\subsection{Diagnostic Statistics}

After computing the log-ratios, the code logs summary statistics to flag potential issues:

\begin{lstlisting}
for name, arr in [("log_r_lin_pos",  log_r_lin_pos), ...]:
    log.info(f"{name}: mean={arr.mean():.3f}, std={arr.std():.3f}")
\end{lstlisting}

What to look for:
\begin{itemize}
  \item \textbf{Mean $\approx 0$}: A zero mean is expected if the auxiliary flow is similar to the SM.
        Large negative mean (e.g.\ $-0.5$) means the auxiliary flow systematically under-rates the
        SM sample---its training distribution is concentrated in a \emph{different} region of phase
        space than the SM samples used here.
  \item \textbf{Small std ($\lesssim 0.5$)}: A small standard deviation indicates that the two flows
        broadly agree on event probability across the sample.  Large std (e.g.\ $> 1.0$) indicates
        strong local variation---some events are much more likely under one flow than the other.  This
        is physically expected for components with very different phase-space support (e.g.\
        \texttt{linear\_neg} trains on only the 16.8\% of events with negative linear weights, so
        its log-ratios naturally have a larger spread).
\end{itemize}

---

\section{The Combination Formula}
\label{sec:combo}

For any desired Wilson coefficient $c_{\text{Ht}}$, the unnormalized weight for event $i$ is:

\begin{equation}
\tilde{w}_i(c) = Z_{\text{sm}} + c \left(Z_{\text{lin}}^+ e^{\log r_{\text{lin},i}^+} - Z_{\text{lin}}^- e^{\log r_{\text{lin},i}^-}\right) 
                 + c^2 \left(Z_{\text{quad}}^+ e^{\log r_{\text{quad},i}^+} - Z_{\text{quad}}^- e^{\log r_{\text{quad},i}^-}\right)
\end{equation}

The normalized weight (for density estimation equivalence at inference time) is:

\begin{equation}
w_i(c) = \frac{\tilde{w}_i(c)}{\sum_j \tilde{w}_j(c)}
\end{equation}

\subsection{Implementation}

\begin{lstlisting}
def get_weights(c, Z, log_r_lin_pos, log_r_lin_neg, 
                log_r_quad_pos, log_r_quad_neg):
    """Compute normalised importance weights for c_Ht."""
    
    unnorm = (Z["sm"]
              + c    * (Z["linear_pos"]    * np.exp(log_r_lin_pos)
                        - Z["linear_neg"]  * np.exp(log_r_lin_neg))
              + c**2 * (Z["quadratic_pos"] * np.exp(log_r_quad_pos)
                        - Z["quadratic_neg"] * np.exp(log_r_quad_neg)))
    
    # Check for destructive interference (negative weights)
    n_neg = (unnorm < 0).sum()
    if n_neg > 0:
        log.info(f"  c={c:+.1f}: {n_neg} events have negative weights")
    
    return unnorm / unnorm.sum()
\end{lstlisting}

\subsection{Interpretation of Negative Weights}

A negative weight $w_i(c) < 0$ at a particular point in phase space reflects \textbf{destructive interference} 
between the SM and EFT amplitudes at that location. This is physically meaningful and expected to occur at 
large Wilson coefficient values. The effective sample size:

\begin{equation}
N_{\text{eff}} = \frac{\left|\sum_i w_i\right|}{\sum_i |w_i|}
\end{equation}

quantifies the fraction of weight budget that survives after positive and negative contributions cancel. 
A sharp drop in $N_{\text{eff}}$ signals that the EFT point is approaching the validity range boundary.

---

\section{Training Pipeline}

\subsection{Submission with \texttt{train\_all\_weights.py}}

\begin{lstlisting}
ALL_WEIGHT_TYPES = ["sm", "linear_pos", "linear_neg", 
                    "quadratic_pos", "quadratic_neg"]

for weight_type in ALL_WEIGHT_TYPES:
    job_cmd = (f"python -m ml.custom.ttz.main_flows "
               f"data_config.weight_type={weight_type} "
               f"data_config.load_weights=True "
               f"data_config.use_weights=True "
               f"training_config.use_sm_weights=True")
    submit_to_condor(job_cmd, weight_type)
\end{lstlisting}

Each weight type is submitted as a separate Condor job. Models are registered in MLflow with names like:
\begin{center}
\texttt{MAFMADEMOG\_flow\_model\_gauss\_rank\_<date>\_cHt5\_linear\_pos\_nall}
\end{center}

\subsection{Loss Function}

The training objective is standard negative log-likelihood with per-batch SNIS normalization:

\begin{equation}
\mathcal{L}_{\text{batch}} = \frac{1}{N_b} \sum_{i \in \text{batch}} \frac{w_i}{\frac{1}{N_b}\sum_j |w_j|} \cdot 
                           \left(-\log_{\det}\left(\frac{\partial f^{-1}}{\partial x}\right) - \log q(x_i)\right)
\end{equation}

where:
\begin{itemize}
  \item The normalization factor $\frac{1}{N_b}\sum_j |w_j|$ stabilizes the loss scale as weights become large
  \item Per-event weights $w_i$ are strictly positive (sign-flipped in the dataset for \texttt{\_neg} types)
  \item Stratified sampling ensures each batch draws equally from high- and low-weight quantile bins
\end{itemize}

\subsection{Data-Adaptive Model Scaling}

Since different weight components occupy different fractions of the event space, models are scaled 
appropriately to their training set size to avoid overfitting:

\begin{equation}
\begin{array}{lccc}
\text{Weight Type} & \text{Training Events} & \text{\% of Total} & \text{Model Scale} \\
\hline
\texttt{sm} & 3.7M & 100\% & 1.00\times \\
\texttt{linear\_pos} & 3.08M & 83.2\% & 1.00\times \\
\texttt{linear\_neg} & 624k & 16.8\% & 0.67\times \\
\texttt{quadratic\_pos} & 3.71M & 99.96\% & 1.00\times \\
\texttt{quadratic\_neg} & 1.5k & 0.04\% & 0.33\times \text{ (sparse/tiny)}
\end{array}
\end{equation}

The scale factor reduces model capacity (hidden dimensions, number of flows) for small datasets. 
For example, \texttt{quadratic\_neg} trains on only 1,453 events with negligible mean weight, so a smaller model 
($33\%$ of default capacity) prevents overfitting while still capturing the structure.

---

\section{Validation and Diagnostics}

\subsection{Effective Sample Size Curve}

After training converges, \texttt{smeft\_reweighting.py} produces an $N_{\text{eff}}$ vs.\ $c_{\text{Ht}}$ curve 
by evaluating the weight formula over a range of $c$ values. A sharp drop indicates:
\begin{itemize}
  \item Destructive interference dominates at that $c$ value
  \item EFT is approaching the validity/perturbativity limit
  \item Fewer events carry significant statistical weight
\end{itemize}

\subsection{Comparison Plots}

For each $c$ value, per-feature comparison panels show:
\begin{enumerate}
  \item \textbf{Generated (red)}: Samples from SM flow, reweighted to $c_{\text{Ht}}$
  \item \textbf{ROOT file (blue, if available)}: Ground-truth baseline from simulation
  \item \textbf{Ratio (lower panel)}: Generated / ROOT, should be $\approx 1$ across phase space
\end{enumerate}

Close agreement validates that the five flows and combination formula accurately capture 
the SMEFT reweighting.

---

\section{File Structure and Key Functions}

\begin{table}[h!]
\centering
\small
\begin{tabular}{lp{10cm}}
\toprule
\textbf{File} & \textbf{Key Functions/Sections} \\
\midrule
\texttt{ml/custom/ttz/ttz\_dataset.py} & \texttt{ttzDataModule.setup()} — loads data, filters by sign, 
normalizes weights \\
& \texttt{WeightStratifiedSampler} — creates stratified batches \\
\midrule
\texttt{ml/custom/ttz/train\_all\_weights.py} & \texttt{ALL\_WEIGHT\_TYPES} — list of weight types \\
& \texttt{submit\_job()} — submits Condor job \\
\midrule
\texttt{ml/custom/ttz/smeft\_reweighting.py} & \texttt{compute\_z\_normalisations()} — computes $Z_k$ \\
& \texttt{find\_model\_name()} — finds latest MLflow models \\
& \texttt{compute\_log\_ratios()} — computes 4 log-ratio arrays \\
& \texttt{get\_weights()} — combination formula \\
& \texttt{plot\_reweighted\_distributions()} — per-$c$ validation plots \\
& \texttt{plot\_neff\_curve()} — $N_{\text{eff}}$ vs.\ $c$ \\
\bottomrule
\end{tabular}
\end{table}

---

\section{Mathematical Consistency Checks}

\subsection{Decomposition Exactness}

The assertion in \texttt{compute\_z\_normalisations()}:
\begin{equation}
\left| Z_{\text{lin}} - (Z_{\text{lin}}^+ - Z_{\text{lin}}^-) \right| < 10^{-3}
\end{equation}

verifies that the decomposition is exact to machine precision.

\subsection{Weight Positivity in Datasets}

For each \texttt{\_pos}/\texttt{\_neg} dataset split, all training weights satisfy $w_i \geq 0$ and 
$\text{mean}(w_i) = 1$ by construction (global mean normalization).

\subsection{Combination Formula Reconstruction}

At $c = 0$, the combination formula reduces to $w_i(0) = Z_{\text{sm}} / Z_{\text{sm}} = 1$ for all $i$, 
giving a uniform distribution—correct since there is no reweighting.

---

\section{Summary and Advantages}

\subsection{Why Pos/Neg Decomposition Solves the Problem}

\begin{enumerate}
  \item \textbf{Principled}: Exact mathematical decomposition, not a hack
  \item \textbf{Stable}: All training weights are non-negative; objective is well-posed
  \item \textbf{Complete}: No events discarded; all contribute to one of five flows
  \item \textbf{Physically correct}: Signed weights at inference (destructive interference) are preserved
  \item \textbf{Validated}: Combination formula verified to exactly reconstruct original signed distribution
  \item \textbf{Robust}: Works without clamping, absolute values, or lossy rescaling
\end{enumerate}

\subsection{Expected Outcomes}

When training converges:
\begin{itemize}
  \item Linear and quadratic flows learn their respective non-negative weight distributions
  \item Validation loss remains stable (no overflow) throughout training
  \item Effective sample size stays $> 10\%$ for a wide range of $c_{\text{Ht}}$ values
  \item Generated event distributions closely match ROOT file at all tested $c$ values
\end{itemize}

\subsection{Recent Updates (March 2026)}

\begin{itemize}
  \item \textbf{Data-Adaptive Scaling}: Models automatically scaled to dataset size to prevent overfitting on sparse components 
        (\texttt{quadratic\_neg} uses only 33\% of default capacity)
  \item \textbf{Data Source Consistency}: Unified all pipeline components to use \texttt{ttZ\_for\_Melle\_tree.root} as the 
        authoritative source file. Z-factor computation now matches training data exactly.
  \item \textbf{Sample Analyzer Fixes}: Fixed \texttt{analyzer.py} bugs (UnboundLocalError for \texttt{\_pos/\_neg} types, 
        greedy regex swallowing model-name suffixes) ensuring feature comparison plots correctly weight MC data with 
        sign-filtered subsets.
  \item \textbf{Weight Contribution Visualization}: Extended \texttt{plot\_weight\_contributions.py} to decompose weights 
        into four separate pos/neg components, enabling structural analysis of positive vs.\ negative weight regions.
\end{itemize}

\subsection{Future Directions}

Once stable training is confirmed:
\begin{enumerate}
  \item Extend to other SMEFT coupling combinations (e.g., $c_t$, $c_W$)
  \item Study higher-order (cubic, quartic) EFT terms if needed
  \item Produce unweighted "pseudo-data" samples for phenomenological analyses
  \item Quantify systematic uncertainties from flow modelling
\end{enumerate}

\appendix

\section{Quick Reference: Weight Treatment Pipeline}

\begin{enumerate}
  \item \textbf{ROOT file}: Extract $w_{\text{sm}}, w_{\text{lin}}, w_{\text{quad}}$ (signed)
  \item \textbf{Dataset.setup()}: Filter by sign $\to$ 5 datasets, each with $w_i > 0$
  \item \textbf{DataLoader}: Use stratified sampler on both train/val $\to$ balanced batches
  \item \textbf{Training}: SNIS normalization $w_i / \text{mean}(|w|)$ stabilizes loss
  \item \textbf{MLflow}: Register 5 models, save checkpoints
  \item \textbf{smeft\_reweighting.py}: Load 5 models, compute 4 log-ratio arrays
  \item \textbf{Combination}: For any $c$, reconstruct $w_i(c)$ using formula in Section \ref{sec:combo}
  \item \textbf{Validation}: Compare reweighted distributions to ROOT file
\end{enumerate}

\end{document}
